% !TEX root = ../../现代密码学简介.tex
\chapter{哈希算法}
\section{消息认证}
我们在进行安全通信的过程中，有许多需要关注的问题。比如说：
\begin{itemize}
	\item 接收者如何确认发送者的身份
	\item 接收者如何确认消息没有被篡改
\end{itemize}

这两个问题在密码学中被称作消息认证(Message Authentication)问题。\par
在上一章讨论的安全协议中，我们可以利用一些复杂的安全协议，达到甚至在没有第三方的情况下，仍然能确认发送者的身份。但是，我们到目前为止却并没有找到能保证消息没有被篡改的方法。\par
因此，我们首先解决第二个问题：如何确认消息没有被篡改。
\section{定义}
在讨论哈希函数之前，我们需要充分讨论一下目前需要完成的目标：接收者如何确认消息没有被篡改。如果一个算法要实现这个目标，那么这个算法必然是接收我们要传递的消息$m$，并产生一个输出$h$。发送者利用这个算法，将消息$m$以及输出$h$都发送给接收者，而接收者再次利用这个算法处理收到的消息$m'$，通过对比输出$h'$与发送者发送的输出$h$，判断消息$m$是否得到篡改。那么，这一算法最重要的一点是什么呢？显然是尽量避免两个不同的消息$m_1\neq m_2$, 其通过这个算法产生的输出$h_1$和$h_2$相同。如果产生这种情况，我们则称该算法拥有一个\textbf{碰撞}。更进一步地，从抵抗攻击的角度来看，这种算法应该使敌手在截获消息经过该算法的输出$h$以后，不能或者难以生成一个篡改的消息$m'$, 使得其经过该算法的输出$h'=h$. 在最理想的状态下，实现这一目标最好的方法就是，从$h$中敌手无法获得任何有关$m$的信息。\par
根据以上的讨论，这个算法$H(x)$需要满足的性质有：
\begin{itemize}
	\item 可以适用于任意长度的消息$x$
	\item 算法的输出$H(x)$是固定长的
	\item 由输入$x$求得输出$H(x)$的过程较容易，可用硬件和软件实现
	\item 由输出$H(x)$反求得输入$x$的过程是不可行的，或者是难以计算的
	\item 给定$x$, 难以求得$y$, 使得$H(x)=H(y)$. 这一点被称作\textbf{弱碰撞抵抗}
	\item 难以求得任何一对$\pth{x, y}$, 使得$H(x)=H(y)$. 这一点被称作\textbf{强碰撞抵抗}
\end{itemize}

对于第二点，即算法的输出是固定长的。这是由于上述讨论中我们提到，从输出中敌手无法获得输入的任何信息，这里甚至包括输入的消息的长度。不仅如此，在实际操作中，由于“由输入求得输出的过程较容易”这一要求，如果输出是固定长的，那么程序中内存分配等问题就相对较容易。因此，现在通行的算法的输出都是固定长度的。同时，我们也应该了解到，这个“固定长”的长度一般是小于消息的长度的。\par
用于确认消息没有篡改的算法被称为哈希函数(Hash fucntion). 又称为散列函数。其接受一个消息作为输入，产生的输出被称为该消息的摘要(digest).\par
一个优秀的哈希函数，应该能做到，对于输入消息的任何一点修改，都会导致输出的巨大改变。下面展示一个例子：\par
如图是东南大学的校徽，我们对其左上角第一个像素进行篡改。
\begin{figure}[H]
\centering
\begin{minipage}[t]{0.49\textwidth}
\centering
\includegraphics[width=10em]{chapters/chapter_6/MD5_unchanged.jpg}
\caption{未篡改的原消息}
\end{minipage}
\begin{minipage}[t]{0.49\textwidth}
\centering
\includegraphics[width=10em]{chapters/chapter_6/MD5_changed.jpg}
\caption{篡改后的消息}
\end{minipage}
\end{figure}

我们利用MD5哈希算法求其摘要：
\begin{itemize}
	\item 未篡改的原消息\par
	\verb`d3f83ad3cee8eb1c407955509b9f08a0`
	\item 篡改后的消息\par
	\verb`d5a15358b5c875e89600b0b53510f5c7`
\end{itemize}

由此可见，即使对原消息进行一丝一毫的修改，都会产生摘要值的巨大改变。\par
除了上述功能之外，由于其输出的长度是固定的，因此，哈希函数还有一个重要的用途：当我们要处理的数据很大时，将其通过哈希函数，压缩成较小的数据，则更好处理。
\section{应用}
在安全通信中，一般较少会直接用到哈希函数，其通常会与之后提到的MAC技术等相结合使用。但是，在一些特定的情况下，哈希函数也是最好的选择。\par
我们思考几个场景：\par
第一，由于某个文件过大，所以传递过程缓慢。但是，在截止日期之前，发送者一定要让接收者了解到该文件已经完成，并且接收者在实际收到文件后，需要确认其最后一次修改是在截止日期前。这个情况虽然少见，但也是确实存在的，比如说数学建模比赛，最后的论文可以在截止日期后提交。解决这一问题的方法就是发送者先对文件求一个哈希，将得到的摘要值发给接收者。等文件真正发送给接收者后，接收者再对文件求哈希，如果两个摘要值相等，就说明在截止日期前文件已经完成修改。\par
第二，在聊天软件中传输文件。现在的机制一般是发送者将发送的文件上传到聊天软件的服务器中，接收者从服务器中下载。那么，为了节省服务器资源，聊天软件设计者一般会在发送者上传文件的时候，对文件求哈希。将摘要值与服务器中已有的文件作比较，如果相同则无需再次上传，接收者可以直接下载对应的文件。\par
下面，我们将具体介绍哈希函数的实现方法
\section{Merkle-Damg\aa rd结构}
在密码学中，许多哈希函数都基于一个类似的结构——Merkle-Damg\aa rd结构。首先，我们先介绍一下该结构的两个主要部分。
\subsection{服从MD结构的填充函数}
对于正整数$b$来说，如果将所有长度为$b$的整数倍的二进制串构成的集合记作$B$, 将所有长度小于$2^b$的二进制串构成的集合记作$U$. 同时，对于长度为$b$的整数倍的二进制串$M$，将其从头开始分割为长度均为$b$的二进制子串，那么称每个子串为$M$的一个块。\par
那么对于非空集合$D\subset U$, 以及函数$\mathrm{Pad}:D\to B$, 称满足以下条件的函数$\mathrm{Pad}$为一个相对于$b$的服从MD结构的填充函数：
\begin{itemize}
	\item 对于任意$M\in D$, $M$是二进制串$\mathrm{Pad}\pth{M}$的前缀
	\item 对于任意$M_1, M_2\in D$, 如果$M_1$和$M_2$的长度相同，那么$\mathrm{Pad}\pth{M_1}$和$\mathrm{Pad}\pth{M_2}$的长度相同
	\item 对于任意$M_1, M_2\in D$, 如果$M_1$和$M_2$的长度不同, 那么$\mathrm{Pad}\pth{M_1}$和$\mathrm{Pad}\pth{M_2}$的最后一个块不同
\end{itemize}

对于任意正整数$b$, 我们可以轻易地构造出服从MD结构的填充函数：对于长度为$k$的二进制串$M$, 如果$k\bmod b\neq b-1$, 那么将二进制串$k$后不停地填充$0$, 直到新串的长度模$b$为$b-1$. 然后，将最初串$M$的长度$k$的二进制表示的最后一位填充在最后，形成一个长度为$b$的整数倍的串。可以轻松地证明，这种填充方法，是服从MD结构的。\par
在上述构造方法中，我们提到，在新串的末尾填充最初串的长度，这一方法被称为Merkle-Damg\aa rd强化。
\subsection{单向压缩函数}
所谓单向压缩函数，其类似于输入是固定长度的哈希函数。对于单向压缩函数$h$, 其接受两个参数，一个为长度为$b$的二进制串，一个为长度为$v$的二进制串，并输出一个长度为$b$的二进制串。这可以看成一个输入长度为$b+v$, 输出长度为$b$的哈希函数。类似于哈希函数，单向压缩函数也有碰撞抵抗等要求。\par
利用单向压缩函数，我们可以把设计接收为任意长度的哈希算法归结到设计接收为固定长度的单向压缩函数上来。那么，单向压缩函数又可以怎么设计呢？\par
而在现行的哈希算法中，单向压缩函数通常利用一个分组密码，使用Davies-Meyer结构构造。具体来说，如果单向压缩函数$h(x, y)$接受长度为$b$比特的$x$和长度为$v$比特的$y$, 输出长度为$b$比特。那么，我们可以使用一个密钥长度为$v$比特，输入长度为$b$比特，输出长度为$b$比特的分组密码加密算法$\E{k}{m}$. Davies-Meyer结构就是
\begin{equation}
h\pth{x, y}=\E{y}{x}\oplus x
\end{equation}

利用Davies-Meyer结构，我们又可以把设计单向压缩函数归结到设计一个抵抗碰撞的分组密码上来。
\subsection{算法}
要实现服从Merkle-Damg\aa rd结构的输出为$b$比特的哈希函数$H$，需要两个正整数$b, v$. 一个相对于$v$的使用了Merkle-Damg\aa rd强化的服从MD结构的填充函数$\mathrm{Pad}$, 一个接受参数为分别为长度为$b$和$v$的二进制串的单向压缩函数$h$, 以及一个长度为$b$的初始值$\mathrm{IV}$.\par
那么，Merkle-Damg\aa rd结构的算法为：
\begin{enumerate}
	\item 将输入的二进制串$M$用$\mathrm{Pad}$函数填充为长度为$nb$比特的二进制串$M'=\mathrm{Pad}\pth{M}$
	\item 将$M'$从头开始分割为$n$个长度为$b$比特的子串$M_1, M_2, \ldots, M_n$
	\item 对于每个子串$M_i$, 计算$H_i=h\pth{M_i, H_{i-1}}$, 其中$H_0=\mathrm{IV}$
	\item $H$的输出即为$M_n$
\end{enumerate}

同时，Merkle和Damg\aa rd分别独立证明出来，如果单向压缩函数$h$能抵抗碰撞，那么$H$也能抵抗碰撞。\par
接下来，介绍两个基于Merkle-Damg\aa rd结构的哈希算法：MD5和SHA-1
\section{MD5}
MD5哈希算法的输出为128比特的二进制串，其$b=128, v=512$. 然而，MD5算法并不安全。在2012年，MD5算法被卡耐基梅隆大学软件工程研究所宣布能被密码学方法破解，不适合将来的使用。但是，如果只是简单的摘要对比，对安全性并不高的话，MD5哈希算法也是一个比较简便的方法。\par
在MD5算法中，二进制串既要以比特串的形式进行异或等逻辑运算，也要以二进制数的形式进行加减等算术运算。因此，在MD5算法中涉及到了数据的存储问题。由于是在计算机上实现的，而计算机中一个数据单元是1字节，也就是8比特。因此，数据存储的问题变成了字节存储顺序的问题。比如说，以十六进制表示的二进制串$1234$究竟表示的是$1\times 16^3+2\times 16^2+3\times 16+4$还是$4\times 16^3+3\times 16^2+2\times 16+1$呢？前者被称为小端法，后者被称为大端法（但无论是小端法还是大端法，每个字节内部比特的存储顺序都是小端法）。而MD5算法中使用的是大端法。但是，现代计算机以及本算法使用的\verb`BigInteger`类均采用小端法处理二进制串。因此，在C++程序实现的时候，每次讲二进制串转化为数字，以及数字转化为二进制串时，都要调整端序。\par
虽然事实上，端序并不会影响加密性，但是，作为一个通行的标准，MD5码仍然需要我们使用大端法。\par
MD5算法包括了四个步骤：
\begin{enumerate}
	\item 填充比特
	\item 填充长度
	\item 初始化变量
	\item 处理消息
\end{enumerate}

其中，填充比特和填充长度就是之前讲的Merkle-Damg\aa rd结构中的填充函数，而初始化链接变量是为了之后处理消息使用的参数，处理消息则是运用单向压缩函数，并进行输出。
\subsection{填充比特}
假设输入消息的长度为$k$, 那么如果$k\bmod{512}<448$, 那么将其后先填充一个`1', 再不停填充`0'直到新串的长度模$512$等于$448$. 如果$k\bmod{512}>448$, 那么先将其后填充一个`1', 然后再不停填充0，填满一个块以后，再继续填充$448$位。\par
其C程序实现如下：
\begin{prove}
\begin{verbatim}
BigInteger appendPaddingBits(BigInteger plainText)
{
    int spareLength = plainText.getLength() % 512;
    if (spareLength < 448)
    {
        int appendLength = 448 - spareLength;
        bool *appendValue = new bool[appendLength];
        *(appendValue + appendLength - 1) = true;
        int index = appendLength - 2;
        while (index > -1)
        {
            *(appendValue + index) = false;
            index--;
        }
        plainText.append(appendValue, appendLength);
    }
    else if (spareLength > 448)
    {
        int appendLength = 512 - spareLength + 448;
        bool *appendValue = new bool[appendLength];
        *(appendValue + appendLength - 1) = true;
        int index = appendLength - 2;
        while (index > -1)
        {
            *(appendValue + index) = false;
            index--;
        }
        plainText.append(appendValue, appendLength);
    }
    return plainText;
}
\end{verbatim}
\end{prove}
\subsection{填充长度}
记最初消息的长度为$k$. 将其表示成二进制串，如果其长度小于64位，则再其首部填0；如果其长度大于64位，则截取其最后64位。\par
这里由于涉及到数向二进制串的转化，需要调整端序。\par
其C程序实现如下：
\begin{prove}
\begin{verbatim}
BigInteger appendLength(BigInteger appendedText, int textLength)
{
    BigInteger appendBigInt(textLength);
    appendBigInt.limitTo(64);
    BigInteger bitSlice[8];
    for (int byte = 0; byte < 8; byte++)
    {
        bitSlice[byte] = appendBigInt.slice(8 * byte, 
                                            8 * (byte + 1));
        appendedText.append(bitSlice[7 - byte]);
    }
    return appendedText;
}
\end{verbatim}
\end{prove}
\subsection{初始化变量}
MD5算法中，一共要初始化4个32位的链接变量，和64个辅助变量。\par
在MD5最初使用的大端法表示下，MD5使用的4个32位链接变量的值分别是：
\begin{itemize}
	\item \verb`A=0x01234567`
	\item \verb`B=0x89abcdef`
	\item \verb`C=0xfedcba98`
	\item \verb`D=0x76543210`
\end{itemize}

所谓链接变量，将其链接起来，就是一个长为128比特的二进制串。在Merkle-Damg\aa rd结构结构中，这个128比特的二进制串就是作为初始值$\mathrm{IV}$.\par
由于这也涉及到数向二进制串的转化，因此，需要调整端序。\par
\verb`T[64]`是一个辅助数组，\verb`T[i]`的值是$4294967296\abs{\sin\pth{i+1}}$.\par
其C程序实现如下：
\begin{prove}
\begin{verbatim}
BigInteger A, B, C, D;
void initializeBuffer()
{
    A = BigInteger("0x67452301");
    B = BigInteger("0xefcdab89");
    C = BigInteger("0x98badcfe");
    D = BigInteger("0x10325476");


    for (int i = 0; i < 64; i++)
    {
        T[i] = BigInteger((unsigned int)(4294967296 * 
                                            fabs(sin(i + 1))));
        T[i].limitTo(32);
    }
}
\end{verbatim}
\end{prove}
\subsection{处理消息}
在处理消息的过程中，需要使用四个辅助函数\verb`FF`, \verb`GG`, \verb`HH`, \verb`II`. 而这四个辅助函数又是由另外四个辅助函数\verb`F`, \verb`G`, \verb`H`, \verb`I`定义的。因此，我们首先定义八个函数：
\begin{prove}
\begin{verbatim}
BigInteger F(BigInteger X, BigInteger Y, BigInteger Z)
{
    return (X & Y) | (~X & Z);
}

BigInteger G(BigInteger X, BigInteger Y, BigInteger Z)
{
    return (X & Z) | (Y & ~Z);
}

BigInteger H(BigInteger X, BigInteger Y, BigInteger Z)
{
    return (X ^ Y) ^ Z;
}

BigInteger I(BigInteger X, BigInteger Y, BigInteger Z)
{
    return Y ^ (X | ~Z);
}

BigInteger X[16];

void FF(BigInteger &A, BigInteger B, BigInteger C, BigInteger D, 
        int k, int s, int i){
    BigInteger tmp = A + F(B, C, D) + X[k] + T[i];
    tmp.limitTo(32);
    A = B + tmp.rotLeft(s);
    A.limitTo(32);
}

void GG(BigInteger &A, BigInteger B, BigInteger C, BigInteger D, 
        int k, int s, int i)
{
    BigInteger tmp = A + G(B, C, D) + X[k] + T[i];
    tmp.limitTo(32);
    A = B + tmp.rotLeft(s);
    A.limitTo(32);
}

void HH(BigInteger &A, BigInteger B, BigInteger C, BigInteger D, 
        int k, int s, int i)
{
    BigInteger tmp = A + H(B, C, D) + X[k] + T[i];
    tmp.limitTo(32);
    A = B + tmp.rotLeft(s);
    A.limitTo(32);
}

void II(BigInteger &A, BigInteger B, BigInteger C, BigInteger D, 
        int k, int s, int i)
{
    BigInteger tmp = A + I(B, C, D) + X[k] + T[i];
    tmp.limitTo(32);
    A = B + tmp.rotLeft(s);
    A.limitTo(32);
}
\end{verbatim}
\end{prove}

在上述程序中，\verb`X[16]`是在下一环节中赋值的变量，每个元素为32比特长的二进制数。\verb`T[64]`是一个上一环节初始化的辅助数组。\par
下一环节则是正式处理消息。首先，根据Merkle-Damg\aa rd结构，需要将处理好的消息分为512比特的块，然后将一个块再分为16个长度为32比特的二进制串\verb`X[16]`。处理的过程中，由于用到了算术运算，因此，也要调整端序。处理的过程就是一个单向压缩函数，其算法就是反复调用之前的\verb`FF`, \verb`GG`, \verb`HH`, \verb`II`四个函数。\par
最后，输出的时候，就是先调整端序，然后拼接并输出。\par
其C程序实现如下：
\begin{prove}
\begin{verbatim}
BigInteger processMessage(BigInteger message)
{
    int N = message.getLength() / 32;
    for (int i = 0; i < N / 16; i++)
    {
        for (int j = 0; j < 16; j++)
        {
            X[j] = message.slice(32 * (i * 16 + j), 
                                    32 * (i * 16 + j + 1));
            BigInteger bitSlice[4] = {X[j].slice(0, 8), 
                                        X[j].slice(8, 16), 
                                        X[j].slice(16, 24), 
                                        X[j].slice(24, 32)};
            X[j] = bitSlice[3];
            X[j].append(bitSlice[2]);
            X[j].append(bitSlice[1]);
            X[j].append(bitSlice[0]);
        }
        
        BigInteger AA = A, BB = B, CC = C, DD = D;
        FF(A, B, C, D, 0, 7, 0);
        FF(D, A, B, C, 1, 12, 1);
        FF(C, D, A, B, 2, 17, 2);
        FF(B, C, D, A, 3, 22, 3);
        FF(A, B, C, D, 4, 7, 4);
        FF(D, A, B, C, 5, 12, 5);
        FF(C, D, A, B, 6, 17, 6);
        FF(B, C, D, A, 7, 22, 7);
        FF(A, B, C, D, 8, 7, 8);
        FF(D, A, B, C, 9, 12, 9);
        FF(C, D, A, B, 10, 17, 10);
        FF(B, C, D, A, 11, 22, 11);
        FF(A, B, C, D, 12, 7, 12);
        FF(D, A, B, C, 13, 12, 13);
        FF(C, D, A, B, 14, 17, 14);
        FF(B, C, D, A, 15, 22, 15);
        
        GG(A, B, C, D, 1, 5, 16);
        GG(D, A, B, C, 6, 9, 17);
        GG(C, D, A, B, 11, 14, 18);
        GG(B, C, D, A, 0, 20, 19);
        GG(A, B, C, D, 5, 5, 20);
        GG(D, A, B, C, 10, 9 ,21);
        GG(C, D, A, B, 15, 14, 22);
        GG(B, C, D, A, 4, 20, 23);
        GG(A, B, C, D, 9, 5, 24);
        GG(D, A, B, C, 14, 9 ,25);
        GG(C, D, A, B, 3, 14, 26);
        GG(B, C, D, A, 8, 20, 27);
        GG(A, B, C, D, 13, 5, 28);
        GG(D, A, B, C, 2, 9 ,29);
        GG(C, D, A, B, 7, 14, 30);
        GG(B, C, D, A, 12, 20, 31);
        
        HH(A, B, C, D, 5, 4, 32);
        HH(D, A, B, C, 8, 11, 33);
        HH(C, D, A, B, 11, 16,34);
        HH(B, C, D, A, 14, 23, 35);
        HH(A, B, C, D, 1, 4, 36);
        HH(D, A, B, C, 4, 11, 37);
        HH(C, D, A, B, 7, 16, 38);
        HH(B, C, D, A, 10, 23, 39);
        HH(A, B, C, D, 13, 4, 40);
        HH(D, A, B, C, 0, 11, 41);
        HH(C, D, A, B, 3, 16, 42);
        HH(B, C, D, A, 6, 23, 43);
        HH(A, B, C, D, 9, 4, 44);
        HH(D, A, B, C, 12, 11, 45);
        HH(C, D, A, B, 15, 16, 46);
        HH(B, C, D, A, 2, 23, 47);
        
        II(A, B, C, D, 0, 6, 48);
        II(D, A, B, C, 7, 10, 49);
        II(C, D, A, B, 14, 15, 50);
        II(B, C, D, A, 5, 21, 51);
        II(A, B, C, D, 12, 6, 52);
        II(D, A, B, C, 3, 10, 53);
        II(C, D, A, B, 10, 15, 54);
        II(B, C, D, A, 1, 21, 55);
        II(A, B, C, D, 8, 6, 56);
        II(D, A, B, C, 15, 10, 57);
        II(C, D, A, B, 6, 15, 58);
        II(B, C, D, A, 13, 21, 59);
        II(A, B, C, D, 4, 6, 60);
        II(D, A, B, C, 11, 10, 61);
        II(C, D, A, B, 2, 15, 62);
        II(B, C, D, A, 9, 21, 63);
        
        A = A + AA;
        A.limitTo(32);
        
        B = B + BB;
        B.limitTo(32);
        
        C = C + CC;
        C.limitTo(32);
        
        D = D + DD;
        D.limitTo(32);
    }
    
    BigInteger output = A.slice(24, 32);
    output.append(A.slice(16, 24));
    output.append(A.slice(8, 16));
    output.append(A.slice(0, 8));
    output.append(B.slice(24, 32));
    output.append(B.slice(16, 24));
    output.append(B.slice(8, 16));
    output.append(B.slice(0, 8));
    output.append(C.slice(24, 32));
    output.append(C.slice(16, 24));
    output.append(C.slice(8, 16));
    output.append(C.slice(0, 8));
    output.append(D.slice(24, 32));
    output.append(D.slice(16, 24));
    output.append(D.slice(8, 16));
    output.append(D.slice(0, 8));
    return output;
}
\end{verbatim}
\end{prove}
\subsection{总过程}
基于上述函数，最终总的MD5哈希算法的过程为：
\begin{prove}
\begin{verbatim}
BigInteger MD5_hash(BigInteger plainText)
{
    BigInteger appendedText = appendPaddingBits(plainText);
    BigInteger message = appendLength(appendedText, 
                                        plainText.getLength());
    initializeBuffer();
    return processMessage(message);
}
\end{verbatim}
\end{prove}
\section{SHA-256}
自1993年开始，美国国家标准技术研究所陆续推出了4组哈希算法，称为安全哈希算法(Secure Hash Algorithm, SHA), 分别为SHA-0, SHA-1, SHA-2, SHA-3.\par
SHA-1已经不再视为可抵御有充足资金、充足计算资源的攻击者。2005年，密码分析人员发现了对SHA-1的有效攻击方法，这表明该算法可能不够安全，不能继续使用，自2010年以来，许多组织建议用SHA-2或SHA-3来替换SHA-1。Microsoft、Google以及Mozilla都宣布，它们旗下的浏览器将在2017年前停止接受使用SHA-1算法签名的SSL证书。因此，这里主要介绍SHA-2中最常用的算法：SHA-256哈希算法。\par
SHA-256也采用的是Merkle-Damg\aa rd结构，其中$b=256, v=512$. 也就是说，SHA-256产生的摘要长度为256比特。此外，与MD5不同的是，SHA-256采用的是大端法存储数据。因此，在实际实现的过程中无需调整端序。\par
类似MD5, SHA-256包括四个步骤：
\begin{enumerate}
    \item 填充比特
    \item 填充长度
    \item 初始化变量
    \item 处理消息
\end{enumerate}
\subsection{填充比特}
SHA-256算法的填充比特步骤与MD5算法一模一样，这里不再赘述。
\subsection{填充长度}
SHA-256算法的填充长度步骤与MD5算法一模一样，但由于SHA-256采用的是大端法，故不需要调整端序。\par
其C程序实现如下：
\begin{prove}
\begin{verbatim}
BigInteger appendLength(BigInteger appendedText, int textLength)
{
    BigInteger appendBigInt(textLength);
    appendBigInt.limitTo(64);
    appendedText.append(appendBigInt);
    return appendedText;
}
\end{verbatim}
\end{prove}
\subsection{初始化变量}
SHA-256算法需要初始化的变量有64个32位辅助变量\verb`K[64]`和8个32位链接变量\verb`H[8]`. 将链接变量链接在一起的256位二进制串就是Merkle-Damg\aa rd结构的初始值$\mathrm{IV}$.\par
其C程序实现如下：
\begin{prove}
\begin{verbatim}
BigInteger K[64];
BigInteger H[8];

void initializeBuffer()
{
    K[0] = BigInteger("0x428a2f98");
    K[1] = BigInteger("0x71374491");
    K[2] = BigInteger("0xb5c0fbcf");
    K[3] = BigInteger("0xe9b5dba5");
    K[4] = BigInteger("0x3956c25b");
    K[5] = BigInteger("0x59f111f1");
    K[6] = BigInteger("0x923f82a4");
    K[7] = BigInteger("0xab1c5ed5");
    K[8] = BigInteger("0xd807aa98");
    K[9] = BigInteger("0x12835b01");
    K[10] = BigInteger("0x243185be");
    K[11] = BigInteger("0x550c7dc3");
    K[12] = BigInteger("0x72be5d74");
    K[13] = BigInteger("0x80deb1fe");
    K[14] = BigInteger("0x9bdc06a7");
    K[15] = BigInteger("0xc19bf174");
    K[16] = BigInteger("0xe49b69c1");
    K[17] = BigInteger("0xefbe4786");
    K[18] = BigInteger("0x0fc19dc6");
    K[19] = BigInteger("0x240ca1cc");
    K[20] = BigInteger("0x2de92c6f");
    K[21] = BigInteger("0x4a7484aa");
    K[22] = BigInteger("0x5cb0a9dc");
    K[23] = BigInteger("0x76f988da");
    K[24] = BigInteger("0x983e5152");
    K[25] = BigInteger("0xa831c66d");
    K[26] = BigInteger("0xb00327c8");
    K[27] = BigInteger("0xbf597fc7");
    K[28] = BigInteger("0xc6e00bf3");
    K[29] = BigInteger("0xd5a79147");
    K[30] = BigInteger("0x06ca6351");
    K[31] = BigInteger("0x14292967");
    K[32] = BigInteger("0x27b70a85");
    K[33] = BigInteger("0x2e1b2138");
    K[34] = BigInteger("0x4d2c6dfc");
    K[35] = BigInteger("0x53380d13");
    K[36] = BigInteger("0x650a7354");
    K[37] = BigInteger("0x766a0abb");
    K[38] = BigInteger("0x81c2c92e");
    K[39] = BigInteger("0x92722c85");
    K[40] = BigInteger("0xa2bfe8a1");
    K[41] = BigInteger("0xa81a664b");
    K[42] = BigInteger("0xc24b8b70");
    K[43] = BigInteger("0xc76c51a3");
    K[44] = BigInteger("0xd192e819");
    K[45] = BigInteger("0xd6990624");
    K[46] = BigInteger("0xf40e3585");
    K[47] = BigInteger("0x106aa070");
    K[48] = BigInteger("0x19a4c116");
    K[49] = BigInteger("0x1e376c08");
    K[50] = BigInteger("0x2748774c");
    K[51] = BigInteger("0x34b0bcb5");
    K[52] = BigInteger("0x391c0cb3");
    K[53] = BigInteger("0x4ed8aa4a");
    K[54] = BigInteger("0x5b9cca4f");
    K[55] = BigInteger("0x682e6ff3");
    K[56] = BigInteger("0x748f82ee");
    K[57] = BigInteger("0x78a5636f");
    K[58] = BigInteger("0x84c87814");
    K[59] = BigInteger("0x8cc70208");
    K[60] = BigInteger("0x90befffa");
    K[61] = BigInteger("0xa4506ceb");
    K[62] = BigInteger("0xbef9a3f7");
    K[63] = BigInteger("0xc67178f2");
    
    H[0] = BigInteger("0x6a09e667");
    H[1] = BigInteger("0xbb67ae85");
    H[2] = BigInteger("0x3c6ef372");
    H[3] = BigInteger("0xa54ff53a");
    H[4] = BigInteger("0x510e527f");
    H[5] = BigInteger("0x9b05688c");
    H[6] = BigInteger("0x1f83d9ab");
    H[7] = BigInteger("0x5be0cd19");
}
\end{verbatim}
\end{prove}
\subsection{处理消息}
在处理消息的过程中，需要用到6个辅助函数\verb`CH`, \verb`MAJ`, \verb`BSIG0`, \verb`BSIG1`, \verb`SSIG0`, \verb`SSIG1`. 下面，先定义这6个辅助函数：
\begin{prove}
\begin{verbatim}
BigInteger CH(BigInteger x, BigInteger y, BigInteger z)
{
    return (x & y) ^ (~x & z);
}

BigInteger MAJ(BigInteger x, BigInteger y, BigInteger z)
{
    return (x & y) ^ (x & z) ^ (y & z);
}

BigInteger BSIG0(BigInteger x)
{
    return x.rotRight(2) ^ x.rotRight(13) ^ x.rotRight(22);
}

BigInteger BSIG1(BigInteger x)
{
    return x.rotRight(6) ^ x.rotRight(11) ^ x.rotRight(25);
}

BigInteger SSIG0(BigInteger x)
{
    BigInteger shifted = x >> 3;
    shifted.limitTo(32);
    return x.rotRight(7) ^ x.rotRight(18) ^ shifted;
}

BigInteger SSIG1(BigInteger x)
{
    BigInteger shifted = x >> 10;
    shifted.limitTo(32);
    return x.rotRight(17) ^ x.rotRight(19) ^ shifted;
}
\end{verbatim}
\end{prove}

在处理消息的过程中，首先将消息分为512个比特的块，然后对每个块进行处理。先将每个块分为16个32比特的二进制串，再将其利用辅助函数混合，产生一共64个32比特的二进制串。然后对其交换、逻辑运算，最终产生输出。\par
其C程序实现如下：
\begin{prove}
\begin{verbatim}
BigInteger processMessage(BigInteger message)
{
    int N = message.getLength() / 512;
    for (int i = 0; i < N; i++)
    {
        BigInteger Mi = message.slice(512 * i, 512 * (i + 1));
        BigInteger W[64];
        for (int t = 0; t < 16; t++)
            W[t] = Mi.slice(32 * t, 32 * (t + 1));
        for (int t = 16; t < 64; t++)
        {
            W[t] = SSIG1(W[t - 2]) + W[t - 7] + SSIG0(W[t - 15])
                                                 + W[t - 16];
            W[t].limitTo(32);
        }
        
        BigInteger a = H[0];
        BigInteger b = H[1];
        BigInteger c = H[2];
        BigInteger d = H[3];
        BigInteger e = H[4];
        BigInteger f = H[5];
        BigInteger g = H[6];
        BigInteger h = H[7];
        
        for (int t = 0; t < 64; t++)
        {
            BigInteger T1 = h + BSIG1(e) + CH(e, f, g) + K[t]
                                                         + W[t];
            BigInteger T2 = BSIG0(a) + MAJ(a, b, c);
            h = g;
            g = f;
            f = e;
            e = d + T1;
            e.limitTo(32);
            d = c;
            c = b;
            b = a;
            a = T1 + T2;
            a.limitTo(32);
        }
        
        H[0] += a;
        H[0].limitTo(32);
        H[1] += b;
        H[1].limitTo(32);
        H[2] += c;
        H[2].limitTo(32);
        H[3] += d;
        H[3].limitTo(32);
        H[4] += e;
        H[4].limitTo(32);
        H[5] += f;
        H[5].limitTo(32);
        H[6] += g;
        H[6].limitTo(32);
        H[7] += h;
        H[7].limitTo(32);
    }
    BigInteger output = H[0];
    output.append(H[1]);
    output.append(H[2]);
    output.append(H[3]);
    output.append(H[4]);
    output.append(H[5]);
    output.append(H[6]);
    output.append(H[7]);
    return output;
}
\end{verbatim}
\end{prove}
\subsection{总过程}
基于上述函数，最终总的SHA-256哈希算法的过程为：
\begin{prove}
\begin{verbatim}
BigInteger SHA_256_hash(BigInteger plainText)
{
    BigInteger appendedText = appendPaddingBits(plainText);
    BigInteger message = appendLength(appendedText,
                                         plainText.getLength());
    initializeBuffer();
    return processMessage(message);
}
\end{verbatim}
\end{prove}